\section{Del aprendizaje por imitación al aprendizaje por refuerzo}

\subsection{Espacio del problema y salto de capacidad}

La clase anterior presentó el aprendizaje por imitación: dadas demostraciones expertas, se aprende una política mediante aprendizaje supervisado. Pero los dos cuellos de botella centrales del aprendizaje por imitación siguen presentes: no puede superar a quien demuestra y carece de un mecanismo de práctica que se refuerce a sí mismo. El aprendizaje por refuerzo introduce el concepto de reward, coloca la conducta en el contexto de una decisión a largo plazo y, por eso, puede mejorar de manera gradual mediante una gran cantidad de experiencia.

\begin{knowledgebox}{La diferencia entre Offline y Online}
\begin{itemize}
    \item \textbf{Offline}: solo puede reutilizar datos históricos; es adecuado para escenarios donde no se puede interactuar libremente (simuladores/medicina).
    \item \textbf{Online}: puede recolectar datos nuevos de manera continua y permite que la política explore en el entorno real.
\end{itemize}
El gradiente de política pertenece al RL online: toma la política de comportamiento directamente como variable de optimización.
\end{knowledgebox}

\subsection{El objetivo de optimización del aprendizaje por refuerzo}

El objetivo final en el aprendizaje por refuerzo es maximizar el trajectory reward: $J(\theta) = \mathbb{E}_{\tau \sim p_\theta(\tau)} [\sum_t r_t]$. Ya no intentamos ajustar las acciones expertas, sino que hacemos descenso (o ascenso) de gradiente en el espacio de $\theta$, de modo que la política favorezca cada vez más las trayectorias con reward alto.

\begin{importantbox}{¿Por qué no hacer aprendizaje supervisado directamente?}
\begin{itemize}
    \item Los datos de comportamiento provienen de la distribución de políticas distintas, por lo que no se puede suponer que sean IID;
    \item el reward ofrece una retroalimentación para explorar desde cero, capaz de descubrir políticas que superen al experto;
    \item en el RL, $\pi_\theta$ determina la recolección de datos, de modo que la política y los datos cambian de manera conjunta.
\end{itemize}
\end{importantbox}

\subsection{El flujo básico del RL online}

\begin{enumerate}
    \item Inicializar $\pi_\theta$ (aleatoria, por imitación, heurística).
    \item Ejecutar $\pi_\theta$ para muestrear varias trajectory (collect data).
    \item Estimar el gradiente de política y actualizar los parámetros $\theta$.
    \item Continuar el muestreo con la nueva política, formando un ciclo cerrado.
\end{enumerate}

\subsection{Resumen de la sección}

Esta sección, desde la perspectiva de la forma del problema y el salto de capacidad, enfatiza que la relación entre el gradiente de política y el aprendizaje por imitación radica en la actualización de la distribución de muestreo, el modelado explícito del objetivo de reward y la capacidad de buscar soluciones con reward alto dentro de un espacio de políticas más amplio.

